% ============================================================
% Paper 5
% Spectral Trace Formula and Smoothed Zero Sums:
% A Prime–Zero Duality Framework
%
% Author:   Ulrich Tehrani
% Date:     May 2026
% Version: v1.19
% Zenodo:   https://doi.org/10.5281/zenodo.19508547
% License:  CC BY 4.0
% ============================================================

\documentclass[11pt]{article}
\usepackage{cmap}            % ToUnicode CMaps for §, en/em-dashes etc.
\usepackage[T1]{fontenc}
\emergencystretch=1em           % avoid minor overfull hboxes
\usepackage[utf8]{inputenc}
\usepackage[a4paper,margin=1in]{geometry}
\usepackage{amsmath,amssymb,amsthm,mathtools}
\usepackage{xcolor}
\usepackage{hyperref}
\usepackage{enumitem}
\usepackage{setspace}
\usepackage{booktabs}
\usepackage{array}
\usepackage{graphicx}

\hypersetup{colorlinks=true,
  linkcolor=blue!50!black,
  urlcolor=blue!50!black,
  citecolor=blue!50!black,
  pdftitle={Spectral Trace Formula and Smoothed Zero Sums:
            A Prime--Zero Duality Framework},
  pdfauthor={Ulrich Tehrani},
  pdfsubject={Research Note, Riemann zeta function,
              prime--zero duality},
  pdfkeywords={Riemann zeta function, trace formula,
               Guinand--Weil, prime-zero coupling}}
\setlength{\parskip}{0.5em}
\setlength{\parindent}{0pt}
\onehalfspacing

% --------------------------------------------------------
% Theorem environments (numbered within section)
% --------------------------------------------------------
\newtheorem{theorem}{Theorem}[section]
\newtheorem{proposition}[theorem]{Proposition}
\newtheorem{lemma}[theorem]{Lemma}
\newtheorem{corollary}[theorem]{Corollary}
\newtheorem{definition}[theorem]{Definition}
\newtheorem{assumption}[theorem]{Assumption}
\newtheorem{observation}[theorem]{Observation}
\newtheorem{conjecture}[theorem]{Conjecture}
\newtheorem*{remark}{Remark}
\newtheorem{openproblem}[theorem]{Open Problem}

% Heuristic environment (clearly marked, never used in proofs)
\newenvironment{heuristic}[1]{%
  \par\smallskip\noindent\textit{[Heuristic: #1]}\quad}{%
  \par\smallskip}

% --------------------------------------------------------
% Convention box — clean ruled box, no mdframed
% --------------------------------------------------------
\newenvironment{conventionbox}{%
  \medskip
  \noindent\rule{\linewidth}{0.4pt}\smallskip\par\noindent
}{%
  \par\smallskip\noindent\rule{\linewidth}{0.4pt}
  \medskip
}

% --------------------------------------------------------
% Series-wide macros
% --------------------------------------------------------
\newcommand{\Hstr}{H_{\mathrm{str}}}
\newcommand{\Hnull}{H_{\mathrm{null}}}
\newcommand{\Estr}{E_{\mathrm{str}}}
\newcommand{\Espec}{E_{\mathrm{spec}}}
\newcommand{\etaorig}{\eta_{\mathrm{orig}}}
\newcommand{\etaren}{\eta_{\mathrm{ren}}}
\newcommand{\Tt}{\widetilde{T}}
\newcommand{\Gun}{G^{\mathrm{un}}}
\newcommand{\Hlocal}{H_{\mathrm{local}}}

% Smoothed zero sums and trace objects
\newcommand{\Zt}{\widetilde{Z}}
\newcommand{\DSEL}{D_{\mathrm{SEL}}}

% Constants (always use subscript — never plain C)
\newcommand{\Ceta}{C_{\eta}}
\newcommand{\CT}{C_{T}}

% --------------------------------------------------------
% Title
% --------------------------------------------------------
\title{\vspace{-1.5em}\bfseries Spectral Trace Formula and Smoothed Zero Sums:\\[0.4em]
\large A Prime--Zero Duality Framework}
\author{Ulrich Tehrani}
\date{Research Note \textperiodcentered\ May 2026
      \textperiodcentered\ v1.19}

% ============================================================
\begin{document}
\maketitle
% ============================================================

\begin{abstract}
We study a one-parameter family of finite-rank operators $\Tt(\sigma)$
acting on a finite-dimensional Hilbert space $\Hnull$ over the imaginary
parts of the non-trivial zeros of $\zeta(s)$.
The operators arise from an explicit coupling map
$\Phi(\sigma)\colon\Hstr\to\Hnull$
encoding prime--zero resonance via
$\sin(\sigma\gamma_k\log p)$.
No hypothetical input is used: no Riemann Hypothesis,
no Hilbert--P\'olya postulate, no GUE conjecture.

\textbf{What is proved.}
We establish an exact trace formula
$\mathrm{Tr}(\Tt(\sigma)) = \DSEL - O(\sigma)$
algebraically, for all $\sigma\in\mathbb{R}$, by a two-line
trigonometric computation (Theorem~\ref{thm:trace}).
The decomposition $B = \sum_{p\leq\kappa}(\log p)^2\,\mathrm{Re}(Z_p^{(2)})$
is proved from first principles (Proposition~\ref{prop:decomp}).
These results require no analytic assumptions.

\textbf{What is numerical.}
The limiting energy asymmetry satisfies
numerical extrapolation gives $\etaren(\kappa)\to 0.81$
on the tested grid $\kappa\in\{23,53,101,199,503,1009\}$
[limit not proved, Observation~\ref{obs:etainf}].
Three spectral signatures at $\sigma=\tfrac{1}{2}$ --- a trace spike
of $+47\%$, a dominant-eigenvalue maximum, and a spectral-gap minimum
--- are observed numerically [Numerical, \S\,\ref{sec:ops}].
The $\gamma^2$-weighted curvature-bias quantity satisfies $B = -19342.5 < 0$
at $\kappa=53$, $\varepsilon=0.05$, $N=100$ [Numerical].

\textbf{What remains open.}
The Bias Conjecture --- for each fixed prime $p$,
$Z_p^\infty(\varepsilon)<0$ for all
sufficiently small $\varepsilon>0$ --- is the
central open problem (Conjecture~\ref{conj:bias}).
Its proof, together with Proposition~\ref{prop:sign},
would establish the sign of the unweighted integrated bias
$B_{\mathrm{int}}$.
Closing the curvature-bias argument for $O''(\tfrac{1}{2})=-2B$
additionally requires the Weighted Bias Bridge from $B_{\mathrm{int}}$
to the $\gamma_k^2$-weighted quantity $B$,
or an independent proof that $B<0$; the sign transfer to
$B_{\mathrm{int}}$ is additionally subject to the auxiliary
truncation-error condition of Proposition~\ref{prop:sign}.
The finite-cutoff identity $\etaren(\kappa)=1-m_1(\nu_\kappa^{(w)})$
holds exactly; convergence of the weighted limiting spectral measure
$\nu_\infty^{(w)}$ and identification of its limiting first moment are open
[Open, \S\,\ref{sec:asymp}].

\textbf{Structure.}
The proved and conditional statements in
\S\S\,\ref{sec:ops}--\ref{sec:smoothed}
are non-circular and status-separated.
No use of the Riemann Hypothesis, the GUE conjecture,
the Montgomery pair correlation conjecture, or the Hilbert--P\'olya
postulate is made anywhere in \S\S\,\ref{sec:ops}--\ref{sec:smoothed}.
\end{abstract}

\vspace{2em}

% --------------------------------------------------------
% Notation and Conventions  (template invariant)
% --------------------------------------------------------
\section*{Notation and Conventions}
\label{sec:notation}

\emph{Critical line.}
Throughout the series, the critical line means $\Re(s)=\tfrac12$.
The parameter $\sigma$ denotes the real part of $s$, so the critical
line corresponds to $\sigma=\tfrac12$.
The reference point $\sigma_0=\tfrac12$ is chosen as the object of
study, not assumed to be the unique zero location.

\medskip
\emph{Global parameters.}
Unless otherwise stated, computations use the reference parameters
$\kappa=53$ (prime cutoff, giving $\pi(\kappa)=16$ active primes),
$\varepsilon=0.05$ (Gaussian damping parameter),
$N=100$ (number of Riemann zero ordinates), and
$\sigma_0=\tfrac12$.
The zero ordinates $\gamma_1<\gamma_2<\cdots<\gamma_N$ are taken
as fixed numerical inputs; the last ordinate is $\gamma_{100}=236.524$.

\medskip
\emph{Main objects.}
The coupling map $\Phi(\sigma)\colon\Hstr\to\Hnull$ is defined by
\[
\Phi(\sigma)_{k,p}
  :=e^{-\varepsilon^2\gamma_k^2/2}\sin(\sigma\gamma_k\log p).
\]
The operator family $\Tt(\sigma):=\Phi(\sigma)\circ\Phi(\sigma)^*$
acts on $\Hnull$ and satisfies the trace formula
$\mathrm{Tr}(\Tt(\sigma))=\DSEL-O(\sigma)$, where
$\DSEL:=\tfrac12 A(\varepsilon,N)\cdot\pi(\kappa)=10.985$ is
$\sigma$-independent and $O(\sigma)$ carries the full
$\sigma$-dependence:
\[
O(\sigma):=\tfrac12\sum_{k,p}e^{-\varepsilon^2\gamma_k^2}
\cos(2\sigma\gamma_k\log p).
\]
The curvature-bias quantity is
$B:=\sum_{p\leq\kappa}(\log p)^2\mathrm{Re}(Z_p^{(2)})$, where
the $\gamma_k^2$-weighted phasor sum is
\[
Z_p^{(2)}:=\sum_{k=1}^N \gamma_k^2\,
e^{-\varepsilon^2\gamma_k^2}\,p^{i\gamma_k},
\]
the unweighted phasor sum is
$Z_p:=\sum_{k=1}^N e^{-\varepsilon^2\gamma_k^2}p^{i\gamma_k}$
(Definition~\ref{def:Zp}), and the Guinand--Weil object is
$\widetilde Z_p^{\mathrm{GW}}(\varepsilon):=\sum_\rho h^{\mathrm{even}}_{p,\varepsilon}
\bigl((\rho-\tfrac12)/i\bigr)$.

\medskip
\emph{Weights.}
The energy asymmetry $\etaren$ recalled in \S\,\ref{sec:asymp}
uses the weight vector $c_p = c_p^{\mathrm{ren}} := \sqrt{f_p}$,
where $f_p := 4(\log p)^2(2p-1)/\bigl(p(p-1)^2\bigr)$
(Paper~4's convention), giving $\etaren(53)=0.66927$.
This specialisation differs from Paper~3's $\etaorig$,
which uses $c_p^{\eta}:=\sqrt{V_p(\tfrac{1}{2})}$
and gives the numerically different value $0.69078$;
see Paper~3, Notation and Conventions.

\medskip
\emph{Not to be confused with.}
The trace-formula constant $\DSEL=10.985$ (this paper) is distinct
from the diagonal energy $D\approx 9.470$ of Paper~2; the subscript
$\mathrm{SEL}$ is mandatory.
The weighted sum $Z_p^{(2)}$ (Definition~\ref{def:Zp2}, carries
$\gamma_k^2$, enters the curvature decomposition of~$B$) is distinct
from the unweighted sum $Z_p$ (Definition~\ref{def:Zp}, enters the
Bias Conjecture).
The finite truncation $Z_p(\varepsilon)$ (Definition~\ref{def:Zp})
is the half-sum over positive ordinates $k=1,\ldots,N$.
Its real part $Z_{p,N}^+(\varepsilon) := \mathrm{Re}\,Z_p(\varepsilon)$
is the object used in Proposition~\ref{prop:sign}.
The infinite analogue $Z_p^\infty(\varepsilon)
:= 2\sum_{k=1}^\infty e^{-\varepsilon^2\gamma_k^2}\cos(\gamma_k\log p)$
is the object of the Bias Conjecture;
$Z_{p,N}^+ = \tfrac{1}{2}(Z_p^\infty + R_{p,N}^\infty)$
with tail $|R_{p,N}^\infty|\lesssim 10^{-61}$ at $N=100$, $\varepsilon=0.05$.
The full Guinand--Weil sum $\widetilde Z_p^{\mathrm{GW}}$ (all zeros of $\zeta$)
provides formal motivation via the explicit formula but is not used
in computations; see Definition~\ref{def:Ztilde}.
The curvature sum $B=-19\,342.5$ (which includes
$\gamma_k^2$ weights) is distinct from the integrated bias
$B_{\mathrm{int}}=-42.21$ (which does not).
The constant $\CT\approx 2\pi(\kappa)$ governing eigenvalue decay
must not be confused with $\Ceta\approx 0.39$ governing the
renormalized spectral radius.

\newpage

% -------------------------------------------------------
\section{Introduction}
\label{sec:intro}

\subsection{Background and Motivation}

This paper is the fifth in a series developing an operator-theoretic
framework for studying the distribution of the non-trivial zeros of the
Riemann zeta function $\zeta(s)$. The series is built around a
finite-dimensional Hilbert-space model in which prime numbers
mediate coupling between zeros, encoded in an explicit operator family.

\textbf{Paper~1}~\cite{Paper1} introduced the local curvature function
$\Hlocal(\sigma,\kappa) :=
     \psi_1(\sigma) + \sum_{p\leq\kappa}V_p(\sigma)$
(see \textbf{Paper~1}, Notation and Conventions),
where each summand $V_p(\sigma)>0$ is a per-prime contribution.
A key finding was the divergence
$\Hlocal(\tfrac{1}{2},\kappa)\sim 2(\log\kappa)^2\to\infty$,
contrasted with convergence $\Hlocal(\sigma,\kappa)\to C(\sigma)<\infty$
for all $\sigma>\tfrac{1}{2}$, identifying the critical line as a phase
boundary observed, not postulated.

\textbf{Paper~2}~\cite{Paper2} connected this divergence to the Weil
explicit functional, constructing admissible test functions
for which the curvature divergence becomes a structural feature of
the Weil energy.

\textbf{Paper~3}~\cite{Paper3} made the geometric structure explicit:
two finite-dimensional Hilbert spaces $\Hstr$ (over primes) and
$\Hnull$ (over zeros) connected by a linear map $\Phi$, with loop
operator $T=\Phi^*\circ\Phi$.
The energy asymmetry $\etaorig>0$ was established numerically for
the Paper~3 weight $c_p^\eta=\sqrt{V_p(1/2)}$.
The convergent renormalised weight scale
$f_p=4(\log p)^2(2p-1)/(p(p-1)^2)$, used later for $\etaren$,
belongs to the Paper~2/Paper~4 convention.

\textbf{Paper~4}~\cite{Paper4} introduced the dual operator
$\Tt=\Phi\circ\Phi^*$, acting on $\Hnull$, and established a
conditional analytic proof that $\etaren>0$ under the remainder
control assumption~(E\textsubscript{rem}); see Paper~4, \S\,5.
Non-vanishing on $\mathrm{Re}(s)=1$ is settled unconditionally by
Hadamard~\cite{Hadamard} and de la Vall\'ee Poussin~\cite{dlVP}
(1896); the genuinely open
condition is the remainder estimate for the Weyl-equidistribution
contribution.

The present paper takes a different route.
Rather than pursuing the equidistribution condition of Paper~4,
we establish an exact trace formula for $\Tt(\sigma)$ --- purely
algebraic, requiring no analytic assumptions --- and introduce
the smoothed zero-sum theory that exposes a structural negative
bias at $\sigma=\tfrac{1}{2}$.
The central open problem is the \textbf{Bias Conjecture}:
for each fixed prime~$p$,
$Z_p^\infty(\varepsilon)<0$
for all sufficiently small~$\varepsilon>0$.
Its proof would complete the $\sigma=\tfrac{1}{2}$ Curvature Bias
argument via the Reduction Lemma, subject to the stated auxiliary
conditions.
\textbf{We do not prove the Bias Conjecture in this paper.}

\textbf{Relation to other approaches.}
Connes and Consani~\cite{ConnesConsani21} construct an archimedean
compression of the scaling action to Sonin's space, establishing a
trace formula for the Weil functional at the archimedean place via an
infinite-dimensional projection-based operator.
The present paper's finite-rank operator
$\Tt(\sigma)=\Phi(\sigma)\Phi(\sigma)^*$ on $\Hnull$ is indexed by
explicit Gaussian-damped sine phasors; it is a Gram-type construction,
not a compression of an ad\`elic action.

Connes, Consani, and Moscovici~\cite{CCM25} construct a self-adjoint
operator as a rank-one perturbation of the scaling spectral triple,
achieving numerical alignment with the first fifty zero ordinates.
The present construction differs structurally: $\Tt=\Phi\Phi^*$ is
a Gram-type operator without spectral postulates and without rank-one
perturbation structure. The zero ordinates $\gamma_k$ are external
numerical inputs, not spectrum.

Yakaboylu~\cite{Yakaboylu26} constructs a non-symmetric operator $R$
on a Hilbert space whose spectrum is determined by the zeros of
$\zeta(s)$ via Bombieri's quadratic functional, invoking the
Hilbert--P\'olya pathway.
The present paper does not postulate any operator whose spectrum
equals the zero ordinates; $\Tt$ is explicitly falsified as a
Hilbert--P\'olya operator in \S\,\ref{ssec:spectral}.

\subsection{Main Results}

The paper contains three main components.

The first is Theorem~\ref{thm:trace}, an exact trace formula
$\mathrm{Tr}(\Tt(\sigma)) = \DSEL - O(\sigma)$,
proved algebraically for all $\sigma\in\mathbb{R}$.
Here $\DSEL$ is an explicit $\sigma$-independent constant and
$O(\sigma)$ is a finite oscillatory sum over primes and zeros;
the entire $\sigma$-dependence of the trace resides in $O(\sigma)$.

The second is a numerical observation (Observation~\ref{obs:etainf})
that $\etaren(\kappa)\to 0.81$ on the tested grid,
motivating a weighted limiting spectral measure $\nu_\infty^{(w)}$
with $\eta_\infty = 1-m_1(\nu_\infty^{(w)})$
(the finite-cutoff identity uses $\nu_\kappa^{(w)}$, not the
unweighted empirical measure $\nu_\kappa$; convergence open).
The numerical content is robust for $\kappa\leq 1009$;
the convergence and tightness of $\nu_\kappa^{(w)}$
and the identification of its limiting first moment remain open.

The third is the smoothed zero-sum framework of \S\,\ref{sec:smoothed}.
Proposition~\ref{prop:decomp} decomposes the curvature quantity as
$B = \sum_p(\log p)^2\allowbreak\mathrm{Re}(Z_p^{(2)})$,
and the Structural Reduction (\S\,\ref{ssec:reduction})
formally isolates a candidate main term that is strictly negative
via the Guinand--Weil explicit formula.
The Bias Conjecture (Conjecture~\ref{conj:bias}) is stated as the
central open problem, and Proposition~\ref{prop:sign} records the
conditional sign transfer.

\subsection{Reader Contract}

This paper \emph{proves} unconditionally the exact trace formula
$\mathrm{Tr}(\Tt(\sigma)) = \DSEL - O(\sigma)$ (Theorem~\ref{thm:trace})
and the curvature decomposition
$B = \sum_p(\log p)^2\mathrm{Re}(Z_p^{(2)})$
(Proposition~\ref{prop:decomp}).
Proposition~\ref{prop:sign} is conditional on the Bias
Conjecture and the stated truncation-error condition, and it
transfers the sign only to the unweighted integrated bias
$B_{\mathrm{int}}$, not to the $\gamma_k^2$-weighted curvature
quantity $B$.

This paper \emph{establishes numerically} that
$\etaren(\kappa)$ appears to approach $0.81$ on the tested
grid $\kappa\in\{23,53,101,199,503,1009\}$
[numerical extrapolation; no convergence theorem]
(Observation~\ref{obs:etainf}),
that $\mathrm{Re}(Z_p(\varepsilon))<0$ for $14$ of the $16$
primes $p\leq 53$ at $\varepsilon=0.05$
(exceptions: $p=37$ and $p=53$), and that
$B_{\mathrm{int}} = -42.21 < 0$ at reference parameters.

This paper \emph{leaves open} the Bias Conjecture
(Conjecture~\ref{conj:bias}), the tightness and convergence of
$\nu_\kappa^{(w)}$ and, if a limiting measure $\nu_\infty^{(w)}$
exists, the limiting identity
$\eta_\infty=1-m_1(\nu_\infty^{(w)})$,
as well as analytical control of the $\Gamma$-term in the
Guinand--Weil decomposition.

% -------------------------------------------------------
\section{The Operator Family \texorpdfstring{$\Tt(\sigma)$}{T-tilde(sigma)}}
\label{sec:ops}

\subsection{Hilbert Spaces and Coupling Map}

Let $\kappa>1$ be a prime cutoff and $N\geq 1$ a zero cutoff.
Write $\{p_1<p_2<\cdots<p_P\}$ for the primes up to $\kappa$,
and $\gamma_1<\gamma_2<\cdots<\gamma_N$ for the imaginary parts
of the first $N$ non-trivial zeros of $\zeta(s)$, listed in
increasing order and taken as fixed numerical input parameters.
The reference parameters used throughout are
$\kappa=53$, $\varepsilon=0.05$, $N=100$, $\sigma_0=\tfrac{1}{2}$.

\begin{definition}[Hilbert spaces]
\label{def:spaces}
Set $\Hstr := \ell^2(\{p\text{ prime}: p\leq\kappa\})$
(of dimension~$P$) and $\Hnull := \ell^2(\{\gamma_k: k=1,\ldots,N\})$
(of dimension~$N$), each equipped with the standard inner product.
\end{definition}

\begin{remark}
The values $\gamma_k$ are used as numerical inputs throughout.
No property of their distribution --- in particular, no assertion
that $\mathrm{Re}(\rho_k)=\tfrac{1}{2}$ for the corresponding
zeros $\rho_k$ --- is assumed or used in any proof.
\end{remark}

\begin{definition}[Coupling map and operator family]
\label{def:coupling}
Let $\varepsilon>0$.
Define $\Phi(\sigma)\colon\Hstr\to\Hnull$ by
\[
\Phi(\sigma)_{k,p} := e^{-\varepsilon^2\gamma_k^2/2}\sin(\sigma\gamma_k\log p),
\]
and set
\[
\Tt(\sigma) := \Phi(\sigma)\circ\Phi(\sigma)^*\colon\Hnull\to\Hnull.
\]
Explicitly,
$\Tt(\sigma)_{kl}
  = \sum_{p\leq\kappa} e^{-\varepsilon^2(\gamma_k^2+\gamma_l^2)/2}
    \sin(\sigma\gamma_k\log p)\sin(\sigma\gamma_l\log p)$.
The reference operator is $\Tt := \Tt(\tfrac{1}{2})$.
\end{definition}

The operator $\Tt(\sigma)$ is self-adjoint and positive semi-definite
for all $\sigma$, since it is of the form $\Phi\Phi^*$.
Its rank is at most $P$.
We write $\mu_1(\sigma)\geq\mu_2(\sigma)\geq\cdots\geq\mu_N(\sigma)\geq 0$
for its eigenvalues in non-increasing order.

\subsection{Spectral Properties}
\label{ssec:spectral}

\textbf{$\Tt$ is not a Hilbert--P\'olya operator.}
The Hilbert--P\'olya conjecture postulates a self-adjoint operator
whose eigenvalues are precisely the imaginary parts $\gamma_k$ of the
non-trivial zeros.
$\Tt$ is not such an operator.
Numerical computation ($\kappa=53$, $\varepsilon=0.05$, $N=100$) gives
$\mu_j\sim\CT/\gamma_{k(j)}$ with $\CT\approx 2\pi(\kappa)$
and Pearson correlation $r_1=0.950$,
where $k(j)$ denotes the dominant zero-index of eigenmode~$j$.
On the tested finite list, the observed eigenvalues follow an
inverse-ordinate decay pattern $\mu_j\approx\CT/\gamma_{k(j)}\to0$,
which is structurally incompatible with the Hilbert--P\'olya
requirement $\mu_j=\gamma_j\to\infty$.
$\Tt$ is a resonance operator: $\Tt_{kl}$ measures the prime-mediated
coupling between zeros $\gamma_k$ and $\gamma_l$.
The term \emph{resonance} is used here in the descriptive sense of
oscillatory prime-mediated coupling, not in the spectral sense of
Connes~\cite{Connes}, where noncritical zeros appear as resonances
of an ad\`elic action.

The constant $\CT\approx 2\pi(\kappa)$ governing the eigenvalue decay
is distinct from $\Ceta\approx 0.39$, which governs the growth of the
maximum eigenvalue of the renormalized operator
$\lambda_{\max}(T_{\mathrm{ren}})\approx\Ceta\cdot\pi(\kappa)$
(numerically observed in Paper~4~\cite{Paper4}).
These two constants must not be confused.

\textbf{Three numerical signatures at $\sigma=\tfrac{1}{2}$.}
The following spectral features are observed numerically at
$\sigma=\tfrac{1}{2}$ ($\kappa=53$, $\varepsilon=0.05$, $N=100$).
First, the trace satisfies $\mathrm{Tr}(\Tt(\tfrac{1}{2}))=14.624$,
approximately $47\%$ above the minimum observed on the
sampled local window $\sigma\in[0.3,0.7]$.
Second, the dominant eigenvalue $\mu_1(\sigma)$ attains its maximum
$\mu_1(\tfrac{1}{2})=6.835$ at $\sigma=\tfrac{1}{2}$.
Third, the spectral gap $|\mu_8(\tfrac{1}{2})-\mu_9(\tfrac{1}{2})|=0.0019$
attains its minimum at $\sigma=\tfrac{1}{2}$.
These are numerical observations, not theorems.

% -------------------------------------------------------
\section{The Trace Formula}
\label{sec:trace}

\begin{theorem}[Trace Formula, proved algebraically]
\label{thm:trace}
For all $\sigma\in\mathbb{R}$,
\[
\mathrm{Tr}\bigl(\Tt(\sigma)\bigr)
  = \sum_{j=1}^N \mu_j(\sigma)
  = \DSEL - O(\sigma),
\]
where
\[
\DSEL   := \tfrac{1}{2}\,A(\varepsilon,N)\cdot\pi(\kappa),\qquad
A(\varepsilon,N) := \sum_{k=1}^N e^{-\varepsilon^2\gamma_k^2},
\]
\[
O(\sigma) := \tfrac{1}{2}\sum_{k=1}^N\sum_{p\leq\kappa}
             e^{-\varepsilon^2\gamma_k^2}\cos(2\sigma\gamma_k\log p).
\]
The quantity $\DSEL$ is $\sigma$-independent.
The entire $\sigma$-dependence of the trace resides in $O(\sigma)$.
\end{theorem}

\begin{proof}
From Definition~\ref{def:coupling},
\[
\mathrm{Tr}\bigl(\Tt(\sigma)\bigr)
  = \sum_{k=1}^N \Tt(\sigma)_{kk}
  = \sum_{k=1}^N\sum_{p\leq\kappa}
    e^{-\varepsilon^2\gamma_k^2}\sin^2(\sigma\gamma_k\log p).
\]
Applying $\sin^2(x)=(1-\cos 2x)/2$ to each term:
\[
= \tfrac{1}{2}\sum_{k,p} e^{-\varepsilon^2\gamma_k^2}
  - \tfrac{1}{2}\sum_{k,p} e^{-\varepsilon^2\gamma_k^2}\cos(2\sigma\gamma_k\log p)
  = \DSEL - O(\sigma). \qedhere
\]
\end{proof}

The proof is purely algebraic: no analytic assumptions, no asymptotic
arguments, no input from the Riemann Hypothesis.
The formula is exact for any finite $\kappa$, $N$, $\varepsilon$, $\sigma$.

Numerical verification ($\kappa=53$, $\varepsilon=0.05$, $N=100$,
$\sigma$-grid) gives
$\max_\sigma|\mathrm{Tr}_{\mathrm{spec}}(\sigma)-(\DSEL-O(\sigma))|<10^{-14}$,
confirming the identity at floating-point precision.
The reference values are:
$A(\varepsilon,N)=1.37306$; $\pi(53)=16$;
$\DSEL = 10.9845$; $O(\tfrac{1}{2})=-3.6391$;
$\mathrm{Tr}(\Tt(\tfrac{1}{2}))=14.624$.

\begin{corollary}[Extrema correspondence, proved]
\label{cor:extrema}
The functions $\sigma\mapsto\mathrm{Tr}(\Tt(\sigma))$ and
$\sigma\mapsto -O(\sigma)$ differ by the constant $\DSEL$;
equivalently,
\[
  \mathrm{Tr}(\Tt(\sigma)) + O(\sigma) \equiv \DSEL.
\]
In particular, maxima of $\mathrm{Tr}(\Tt(\sigma))$ correspond to
minima of $O(\sigma)$, and conversely.
\end{corollary}

Numerically, $\sigma=\tfrac{1}{2}$ is a local maximum of
$\mathrm{Tr}(\Tt(\sigma))$, equivalently a local minimum of $O(\sigma)$.
This is a numerical observation.

\begin{remark}[Positioning relative to classical trace formulas]
The decomposition $\mathrm{Tr}(\Tt(\sigma))=\DSEL-O(\sigma)$
has a structural parallel to the Selberg trace formula,
in which a spectral sum equals a sum of geometric data:
$\DSEL$ plays the role of the constant geometric term and
$O(\sigma)$ plays the role of the oscillatory prime-geodesic sum.
This analogy is motivational; no formal connection to the
Selberg trace formula is claimed.

The distinction should be stated precisely.
The classical Selberg trace formula matches Laplace spectra on
hyperbolic surfaces with lengths of closed geodesics; the Arthur
trace formula balances geometric and spectral distributions for
reductive groups; the Connes trace formula~\cite{Connes}
interprets the explicit formula on the ad\`ele-class space with
zeros as a spectral absorption set.
In Burnol's Fourier-analytic approach~\cite{Burnol}, the explicit
formula is recast in operator-theoretic language via a conductor
operator $H=\log|x|+\log|y|$.
By contrast, the identity in Theorem~\ref{thm:trace} is a
finite-rank algebraic trace computation for an explicitly
constructed prime--zero Gram family; it does not arise from
group actions, orbital integrals, or an ad\`elic framework.

Three further points deserve emphasis.
Theorem~\ref{thm:trace} makes no assertion about $\sigma=\tfrac{1}{2}$
specifically; the trace formula is uniform in $\sigma$.
No connection to the Riemann Hypothesis follows from
Theorem~\ref{thm:trace} alone.
Finally, $\mathrm{Tr}(\Tt(\sigma))\neq\Gamma_\kappa(\sigma)$ in general:
the Pearson correlation is $+0.9915$ (numerical), but the ratio is not
constant (coefficient of variation $54.76\%$), owing to a $(\log p)^2$
discrepancy in the definitions, since $\Tt$ uses $\sin(\sigma\gamma_k\log p)$
while $\Gamma_\kappa$ involves $(\log p)^2$ factors explicitly.

The symbol $\DSEL$ carries a mandatory subscript to distinguish it
from the Paper~2 diagonal energy scalar
$D:=\sum_p f_p\approx 9.470$; these two quantities are unrelated.
\end{remark}

% -------------------------------------------------------
\section{Asymptotic Spectral Measure}
\label{sec:asymp}

We recall the energy asymmetry from Papers~3 and~4.
For prime cutoff $\kappa$,
\[
\etaren(\kappa) := 1 - \frac{c^T \Gun c}{\Estr(\kappa)},
\]
where $c=(c_p)_{p\leq\kappa}$ is the canonical weight vector with
$c_p := \sqrt{4(\log p)^2(2p-1)\big/\bigl(p(p-1)^2\bigr)}$;
$\Gun$ is the full unnormalised Gram matrix of the resonance vectors,
$\Gun_{pq}=\langle a_p,a_q\rangle$, so that $c^T\Gun c = \Espec$.
The off-diagonal interaction matrix is $\Gun - D_a$,
where $D_a = \operatorname{diag}(\|a_p\|^2)$;
and $\Estr(\kappa):=\sum_{p\leq\kappa}c_p^2\,\|a_p\|^2$, with
$\|a_p\|^2:=\sum_k e^{-\varepsilon^2\gamma_k^2}\sin^2(\gamma_k\log p)$
(Paper-4 fixed phasor convention, no $\sigma$-factor; this connects with
Paper~4's energy-asymmetry framework and differs from the Paper-5 reference
operator $\Phi(\sigma)$ of \S\,\ref{sec:ops}).
The reference value is $\etaren(53)=0.66927$; on the
tested grid $\kappa\in\{23,53,101,199,503,1009\}$
the values lie in $[0.669, 0.820]$
[numerical; no convergence theorem].
The convergence $\Estr(\kappa)\to\Estr(\infty)$ follows from the bound
$\|a_p\|^2\le A(\varepsilon,N)$ and
\[
  f_p = \frac{4(\log p)^2(2p-1)}{p(p-1)^2}
      = O\!\left(\frac{(\log p)^2}{p^2}\right),
\]
so $\Estr(\kappa)\le A(\varepsilon,N)\sum_{p\le\kappa}f_p$ converges
absolutely as $\kappa\to\infty$.
The value $\Estr(\infty)\approx 6.1$ is numerical at the stated
reference parameters [Numerical].

The finite-grid values of $\etaren(\kappa)$ admit an exact algebraic
description. Assume $\|a_p\|>0$ for all active primes (satisfied
numerically throughout). Set
\[
K_\kappa:=D_a^{-1/2}\Gun D_a^{-1/2},\qquad
v_\kappa:=D_a^{1/2}c/\sqrt{\Estr(\kappa)}.
\]
Then $\|v_\kappa\|=1$ and, writing
$K_\kappa=\sum_j\lambda_j u_j u_j^T$,
\[
\nu_\kappa^{(w)}:=\sum_j|\langle v_\kappa,u_j\rangle|^2\delta_{\lambda_j},
\qquad
\etaren(\kappa)=1-m_1\!\left(\nu_\kappa^{(w)}\right) \quad\text{exactly.}
\]
The unweighted empirical eigenvalue measure
$\nu_\kappa:=(1/N)\sum_j\delta_{\mu_j}$
gives $m_1(\nu_\kappa)=(1/N)\mathrm{Tr}(\Tt(\kappa))$,
which is not identifiable with $c^T\Gun c/\Estr$;
it is included only for comparison.

\begin{observation}[Numerical]
\label{obs:etainf}
For $\kappa\leq 1009$, $\varepsilon=0.05$, $N=100$:
\[
\etaren(\kappa)\to 0.81
\quad\textup{[numerical extrapolation; limit not proved]}.
\]
The open problem is tightness and convergence of $\nu_\kappa^{(w)}$
and identification of its limiting first moment
(Open Problem~\ref{op:measure}).
The extrapolated value near $0.81$ is not identified with
any classical constant ($\eta_\infty\neq\pi^2/6$,
$\eta_\infty\neq\gamma_E$).
\end{observation}

The rank-one update structure in the cutoff variable deserves comment.
In the Paper-5 convention $\sigma=\tfrac12$,
each step $\kappa\to\kappa'$ (adding the next prime $p'$) takes the form
\[
\Tt(\kappa') = \Tt(\kappa) + \Phi(\tfrac12)_{p'}\otimes\Phi(\tfrac12)_{p'}^T,
\]
where $\Phi(\tfrac12)_{k,p'} :=
  e^{-\varepsilon^2\gamma_k^2/2}\sin(\tfrac12\gamma_k\log p')$
is the column of $\Phi(\tfrac12)$ associated with $p'$.
This rank-1 perturbation structure provides a tractable framework for
studying the $\kappa\to\infty$ limit.

\begin{remark}[Relation to the energy-gap structure of Paper~4]
The oscillatory term $O(\sigma)$ at $\sigma=\tfrac12$
shares its trigonometric kernel with the prime-phasor
interference structure of Paper~4~\cite{Paper4}.
Specifically,
\[
\DSEL - O(\tfrac12)
= \sum_{p,k} w_k \sin^2\!\bigl(\tfrac12\gamma_k\log p\bigr)
\]
equals the diagonal of $\widetilde T(\tfrac12)$, while
the energy gap $\Delta$ of Paper~4 measures the
off-diagonal pairwise prime-phasor interference.
Both involve related prime-phasor kernels built
from $\gamma_k\log p$, but with different
normalisations: the present paper uses the
$\sigma$-dependent half-angle trace convention
at $\sigma=\tfrac12$, while Paper~4 uses the
fixed phasor convention without $\sigma$.
\end{remark}

% -------------------------------------------------------
\section{Smoothed Zero Sums and the Bias Conjecture}
\label{sec:smoothed}

\subsection{Objects and Decomposition}
\label{ssec:objects}

We introduce two analytically distinct objects and establish their
relationship to the curvature-bias quantity.

\begin{definition}[Finite truncation]
\label{def:Zp}
For a prime $p\leq\kappa$, set
\[
Z_p := \sum_{k=1}^N e^{-\varepsilon^2\gamma_k^2}\,p^{i\gamma_k}\in\mathbb{C},
\qquad
\mathrm{Re}(Z_p) = \sum_{k=1}^N e^{-\varepsilon^2\gamma_k^2}\cos(\gamma_k\log p).
\]
\end{definition}

\begin{definition}[Guinand--Weil object and Ordinate Proxy]
\label{def:Ztilde}
For a prime $p$ and $\varepsilon>0$, define the \emph{full Guinand--Weil zero sum}
\[
\widetilde Z_p^{\mathrm{GW}}(\varepsilon)
  := \sum_\rho h^{\mathrm{even}}_{p,\varepsilon}
     \!\left(\tfrac{\rho-\frac12}{i}\right),
\]
where the sum ranges over all non-trivial zeros $\rho$ of $\zeta(s)$,
and define the \emph{finite positive-ordinate half-sum}
\[
Z_{p,N}^{+}(\varepsilon)
  := \sum_{k=1}^N e^{-\varepsilon^2\gamma_k^2}
     \cos(\gamma_k\log p)
  = \mathrm{Re}\, Z_p(\varepsilon),
\]
restricted to the first $N$ positive ordinates $\gamma_k$.
The test function and its Fourier transform are
\[
h^{\mathrm{even}}_{p,\varepsilon}(u) := e^{-\varepsilon^2 u^2}\cos(u\log p),
\qquad
\hat{h}_{p,\varepsilon}(x) :=
\frac{\sqrt{\pi}}{2\varepsilon}
\left[
\exp\!\left(-\frac{(2\pi x-\log p)^2}{4\varepsilon^2}\right)
+
\exp\!\left(-\frac{(2\pi x+\log p)^2}{4\varepsilon^2}\right)
\right].
\]
In this paper, all computations use $Z_p(\varepsilon)$
(Definition~\ref{def:Zp}) and the infinite analogue
\[
  Z_p^\infty(\varepsilon)
  := 2\!\sum_{k=1}^{\infty}
     e^{-\varepsilon^2\gamma_k^2}\cos(\gamma_k\log p).
\]
The full Guinand--Weil object $\widetilde Z_p^{\mathrm{GW}}$
serves as formal motivation for the sign of $Z_p^\infty$
via its explicit-formula main term (Structural Reduction below),
but is not used in computations or the Bias Conjecture.
The bare symbol $\widetilde Z_p$ is not used in this paper.
\end{definition}

Equivalently,
\[
Z_p^\infty(\varepsilon)
=2\lim_{N\to\infty}Z_{p,N}^+(\varepsilon),
\]
whenever the Gaussian-damped ordinate sum is taken over
all positive ordinates.
The relationship between $\widetilde Z_p^{\mathrm{GW}}$
and $Z_p^\infty$ requires knowledge of off-critical zeros.
If all non-trivial zeros lie on the critical line, then
by evenness of $h^{\mathrm{even}}_{p,\varepsilon}$,
\[
\widetilde Z_p^{\mathrm{GW}}(\varepsilon)
= Z_p^\infty(\varepsilon)
\]
in the infinite-ordinate limit. Without this assumption,
\[
\widetilde Z_p^{\mathrm{GW}}(\varepsilon)
= Z_p^\infty(\varepsilon) + E_p^{\mathrm{off}}(\varepsilon),
\]
where $E_p^{\mathrm{off}}$ collects contributions from zeros off
the critical line. Controlling this term is
Open Problem~\ref{op:gwbridge}.

The factor two here is a zero-side convention: $Z_{p,N}^+$ sums only
over positive ordinates, while the full even zero sum counts the
paired ordinates $\pm\gamma_k$. Under critical-line localization and
in the infinite-ordinate limit, evenness gives
$\widetilde Z_p^{\mathrm{GW}}=2\lim_N Z_{p,N}^+=Z_p^\infty$.
The Fourier-transform normalization affects the prime-side constants
in the explicit formula; it is recorded separately in the
Guinand--Weil convention below.

\begin{remark}[Relation to Mazhouda's asymptotic sums]
The smoothed zero sum $Z_p^\infty(\varepsilon)$ is formally a
prime-indexed complex specialisation of Mazhouda's Selberg-class
Gaussian sum $\sum_\rho e^{u\rho^2 - v\rho}$
\cite{Mazhouda} under $u = \varepsilon^2$, $v = i\log p$.
Mazhouda establishes asymptotic behaviour for $v \in \mathbb{R}$
(cases $v = u$ and $v = 0$)~\cite{Mazhouda};
the prime-indexed imaginary-$v$ case and the sign of
$\mathrm{Re}(Z_p^\infty)$ are not covered by \cite{Mazhouda}
and constitute the central open problem of this paper.
\end{remark}

The curvature-bias quantity is defined as
\[
B := \sum_{k=1}^N\sum_{p\leq\kappa}
     e^{-\varepsilon^2\gamma_k^2}(\gamma_k\log p)^2\cos(\gamma_k\log p).
\]

\begin{definition}[$\gamma^2$-weighted phasor sum]
\label{def:Zp2}
\[
Z_p^{(2)} := \sum_{k=1}^N \gamma_k^2\,
             e^{-\varepsilon^2\gamma_k^2}\,p^{i\gamma_k}.
\]
\end{definition}

\noindent
The $\gamma_k^2$-weighted phasor sum $Z_p^{(2)}$ carries the
frequency weights needed for the curvature decomposition.
It is distinct from the unweighted sum $Z_p$
(Definition~\ref{def:Zp}), which enters the Bias Conjecture.
The $\gamma_k^2$-weighting in the curvature operator is consistent
with the UV prolate spectral weights studied by Connes and
Moscovici~\cite{ConnesMoscovici22} in the context of prolate-function
approximations to the zero ordinates; the analytic significance of
this weighting remains open.

\begin{proposition}[Decomposition, proved]
\label{prop:decomp}
\[
B = \sum_{p\leq\kappa}(\log p)^2\,\mathrm{Re}(Z_p^{(2)}).
\]
\end{proposition}

\begin{proof}
Expand $\mathrm{Re}(Z_p^{(2)})
= \sum_k \gamma_k^2\,e^{-\varepsilon^2\gamma_k^2}\cos(\gamma_k\log p)$,
multiply by $(\log p)^2$, and exchange the order of summation.
The product $\gamma_k^2\cdot(\log p)^2 = (\gamma_k\log p)^2$
recovers the definition of~$B$.
\end{proof}

\begin{remark}
The unweighted decomposition
$\sum_{p\leq\kappa}(\log p)^2\operatorname{Re}(Z_p) = -42.21$
at reference parameters is a distinct quantity, denoted
$B_{\mathrm{int}}$, the unweighted integrated bias
(ordinate proxy; \S\,\ref{ssec:sign}).
Both $B=-19342.5$ and $B_{\mathrm{int}}=-42.21$ are negative;
the curvature value~$B$ is the one entering the
$\sigma=\tfrac{1}{2}$ selection argument via
$O''(\tfrac12)=-2B$.
\end{remark}

Numerically, $B=-19342.5$ ($\kappa=53$, $\varepsilon=0.05$,
$N=100$, numerical).
Numerically, $\mathrm{Re}(Z_p)<0$ holds for $14$ of the $16$ primes
$p\leq 53$; the two exceptions are $p=37$ and $p=53$
($\varepsilon=0.05$, $N=100$).

\subsection{Structural Reduction}
\label{ssec:reduction}

We apply the Guinand--Weil explicit formula to $h^{\mathrm{even}}_{p,\varepsilon}$
to formally identify the candidate main term of
$\widetilde Z_p^{\mathrm{GW}}(\varepsilon)$.
This formal calculation motivates (but does not prove)
the sign of $Z_p^\infty(\varepsilon)$.

\begin{conventionbox}
\emph{Convention} (Guinand--Weil normalization).
We use the explicit formula in the following form.
For an admissible even test function $h$ with Fourier transform
$\hat{h}(x) := \int_{-\infty}^\infty h(u)e^{-2\pi i ux}\,du$,
the Guinand--Weil formula reads
\[
\sum_\rho h\!\left(\tfrac{\rho-\frac12}{i}\right)
  = -\frac{1}{2\pi}\sum_{n\geq 2}\frac{\Lambda(n)}{\sqrt{n}}
    \bigl[\hat{h}({\textstyle\frac{\log n}{2\pi}})
          + \hat{h}({-\textstyle\frac{\log n}{2\pi}})\bigr]
  + (\Gamma\text{-terms}) + (\text{constant terms}).
\]
The factor $1/(2\pi)$ in front of the prime-power sum reflects
the chosen Fourier-transform normalization.
This convention is consistent with Weil~\cite{Weil} and
Guinand~\cite{Guinand}.
\end{conventionbox}

\textit{Structural Reduction~\ref{ssec:reduction}
(formal calculation: formal isolation of candidate main term; dominance open).}

\medskip

\textit{Step~1 --- Explicit formula template.}
Applying the convention above to $h=h^{\mathrm{even}}_{p,\varepsilon}$:
\[
\widetilde Z_p^{\mathrm{GW}}(\varepsilon)
  = -\frac{1}{2\pi}\sum_{n\geq 2}\frac{\Lambda(n)}{\sqrt{n}}
    \bigl[\hat{h}_{p,\varepsilon}({\textstyle\tfrac{\log n}{2\pi}})
    + \hat{h}_{p,\varepsilon}({-\textstyle\tfrac{\log n}{2\pi}})\bigr]
  + (\Gamma\text{-terms}) + (\text{constant terms}).
\]

\textit{Step~2 --- Fourier transform.}
For $h^{\mathrm{even}}_{p,\varepsilon}(u)=e^{-\varepsilon^2 u^2}\cos(u\log p)$,
the Fourier transform is a sum of two Gaussians:
\[
\widehat h_{p,\varepsilon}(x)
=\frac{\sqrt\pi}{2\varepsilon}
\left[
\exp\!\left(-\frac{(2\pi x-\log p)^2}{4\varepsilon^2}\right)
+
\exp\!\left(-\frac{(2\pi x+\log p)^2}{4\varepsilon^2}\right)
\right].
\]
At the peak $x = (\log p)/(2\pi)$ the two summands contribute
$(\sqrt\pi/2\varepsilon)(1 + e^{-(\log p)^2/\varepsilon^2})$,
where the second term is exponentially suppressed for small $\varepsilon$.

\textit{Step~3 --- Prime-power localization.}
For small $\varepsilon$, the Gaussian $\hat{h}_{p,\varepsilon}$ is
sharply concentrated near $x=(\log p)/(2\pi)$.
In the prime-power sum $\sum_{n\geq 2}(\Lambda(n)/\sqrt{n})\hat{h}(\log n/2\pi)$,
the contribution from $n=p$ (with $\Lambda(p)=\log p$) is the
\textit{formally isolated peak contribution}: other prime powers $n\neq p$ contribute
$\hat{h}_{p,\varepsilon}(\log n / 2\pi)$,
which is exponentially suppressed when
$|\log n - \log p| \gg \varepsilon$;
the quantitative dominance over the local
cluster around $\log p$ is not established
here and constitutes part of the Bias Conjecture.

\textit{Step~4 --- Candidate main term.}
The $n=p$ contribution, using the convention above, is
\begin{align*}
-\frac{1}{2\pi}\frac{\log p}{\sqrt{p}}
\left[\hat{h}_{p,\varepsilon}\!\left(\frac{\log p}{2\pi}\right)
+\hat{h}_{p,\varepsilon}\!\left(-\frac{\log p}{2\pi}\right)\right]
&= -\frac{\log p}{2\sqrt{\pi}\,\varepsilon\sqrt{p}}
  \left(1+e^{-(\log p)^2/\varepsilon^2}\right).
\end{align*}
The dominant contribution uses the leading Gaussian peak;
the factor $e^{-(\log p)^2/\varepsilon^2}$ is
exponentially suppressed for small $\varepsilon$.
This candidate main term is strictly negative for all primes $p>1$.
This is a formal statement about $\widetilde Z_p^{\mathrm{GW}}$;
its transfer to $Z_p^\infty$ is conditional on the GW bridge
(Open Problem~\ref{op:gwbridge}).

\textit{Step~5 --- Remainder.}
All remaining contributions --- prime powers $n=p^k$ with $k\geq 2$,
other primes, the archimedean $\Gamma$-factor functional
attached to $h^{\mathrm{even}}_{p,\varepsilon}$,
constant terms, and trivial-zero contributions --- are grouped into
the $(\Gamma\text{-terms})$ and correction terms.
These are real-valued.
Their combined magnitude relative to the candidate main term
is not controlled quantitatively here.

In summary, the Structural Reduction formally identifies
\[
\widetilde Z_p^{\mathrm{GW}}(\varepsilon)
  = -\frac{\log p}{2\sqrt{\pi}\,\varepsilon\sqrt{p}}
    + (\Gamma\text{-terms}) + (\text{correction terms}).
\]
The candidate main term is strictly negative.
Whether it dominates quantitatively is the content of the
Bias Conjecture (Conjecture~\ref{conj:bias}).

\begin{remark}
The $\Gamma$-terms and correction terms in the Structural Reduction
are real-valued. Their magnitude relative to the candidate main term
$-(\log p)/(2\sqrt{\pi}\,\varepsilon\sqrt{p})$ is not controlled
quantitatively in this paper. Establishing this dominance for the full Guinand--Weil object is the
explicit-formula part of the Bias Conjecture program; transferring it to
$Z_p^\infty$ additionally requires the GW bridge of
Open Problem~\ref{op:gwbridge}.
\end{remark}

The following lemma records a formally trivial but logically
useful reduction.

\begin{lemma}[Reduction Lemma, proved]
\label{lem:reduction}
If $B = -M + E$ with $M>0$ and $|E|<M$, then $B<0$.
\end{lemma}

\begin{proof}
Since $|E|<M$, one has $B=-M+E\leq -M+|E|<0$.
\end{proof}

This lemma isolates the local dominance step: once a negative main
term and a smaller remainder have been identified for the relevant
quantity, its sign is fixed. The remaining analytic content lies in
proving such dominance, and in the separate bridge and stationarity
conditions stated below.

\subsection{Sign Transfer and the Bias Conjecture}
\label{ssec:sign}

\begin{definition}[Finite half-sum and truncation tail]\label{def:trunc}
\begin{align*}
Z_{p,N}^{+}(\varepsilon) &:= \mathrm{Re}\,Z_p(\varepsilon)
  = \sum_{k=1}^N e^{-\varepsilon^2\gamma_k^2}
    \cos(\gamma_k\log p), \\[4pt]
R_{p,N}^{\infty}(\varepsilon)
  &:= 2Z_{p,N}^{+}(\varepsilon) - Z_p^\infty(\varepsilon)
   = -2\sum_{k>N} e^{-\varepsilon^2\gamma_k^2}
     \cos(\gamma_k\log p),
\end{align*}
so that
\[
Z_{p,N}^{+}(\varepsilon)
= \tfrac{1}{2}\bigl(Z_p^\infty(\varepsilon)
  + R_{p,N}^{\infty}(\varepsilon)\bigr).
\]
The tail satisfies
$|R_{p,N}^{\infty}(\varepsilon)| \leq 2\int_{\gamma_N}^{\infty}
e^{-\varepsilon^2 t^2}\,dN(t) \lesssim 10^{-61}$
at $N=100$, $\varepsilon=0.05$, uniformly for $p\leq 53$
[Numerical: Riemann--von Mangoldt density argument,
verified by~\texttt{verify\_paper5.py}].
\end{definition}

\begin{conjecture}[Bias Conjecture, open]
\label{conj:bias}
For each fixed prime $p$, there exists $\varepsilon_0(p)>0$ such that
\[
Z_p^\infty(\varepsilon) < 0
\quad\text{for all } 0 < \varepsilon < \varepsilon_0(p),
\quad
Z_p^\infty(\varepsilon)
  = 2\!\sum_{k=1}^\infty e^{-\varepsilon^2\gamma_k^2}
    \cos(\gamma_k\log p).
\]
\end{conjecture}

\noindent
\textit{Terminological note.}
The word ``bias'' is used here for the sign asymmetry of a smoothed
spectral sum at an individual prime scale~$p$.
This is not a prime-number race in the sense of Chebyshev or
Rubinstein--Sarnak~\cite{RubinsteinSarnak}, which concerns the
comparative distribution of primes in residue classes under GRH.
It is also not a Weil positivity statement or a Li-coefficient
criterion, both of which are global positivity conditions over all
zeros and test functions simultaneously.
Conjecture~\ref{conj:bias} is a pointwise spectral-sign claim for
a single-prime-indexed observable under Gaussian damping.

The evidence for Conjecture~\ref{conj:bias} is threefold.
Structurally, the Structural Reduction of \S\,\ref{ssec:reduction}
formally identifies the candidate main term
$-(\log p)/(2\sqrt{\pi}\,\varepsilon\sqrt{p})<0$.
Numerically, $\mathrm{Re}(Z_p)<0$ holds for $14$ of $16$ primes
$p\leq 53$ [Numerical].
As a structural precedent, Mazhouda~\cite{Mazhouda} derived
Gaussian-smoothed asymptotic formulas for sums over zeros of
$L$-functions in the Selberg class: for $m\geq 2$,
\[
\sum_\rho e^{u\rho^2+(\log m)\rho}
  = -\frac{\Lambda(m)}{\sqrt{4\pi u}} + O_m(1), \quad u\to 0^+,
\]
showing that Gaussian smoothing isolates prime-power contributions
and produces an explicit negative main term.
This provides a structural precedent for the negative bias in smoothed
zero sums at prime scales.
The test functions differ from those used here; a direct application
to $Z_p^\infty(\varepsilon)$ requires showing that the present
kernel falls within Mazhouda's framework.

The conjecture remains open because the quantitative dominance of the
candidate main term over all $\Gamma$-terms and correction terms is
not formally established.

We write
\[
B_{\mathrm{int}}(\kappa,N,\varepsilon)
:=\sum_{p\le\kappa}(\log p)^2 Z_{p,N}^+(\varepsilon)
=\sum_{p\le\kappa}(\log p)^2\operatorname{Re}Z_p(\varepsilon).
\]

\begin{proposition}[Conditional sign transfer]
\label{prop:sign}
Assume:
\begin{enumerate}[label=\textup{(\Alph*)}]
\item $Z_p^\infty(\varepsilon)<0$ for all primes $p\leq\kappa$
  \textup{(Bias Conjecture)};
\item $|R_{p,N}^{\infty}(\varepsilon)|
  < |Z_p^\infty(\varepsilon)|$
  for all primes $p\leq\kappa$
  \textup{(truncation tail strictly dominated by main term)}.
\end{enumerate}
Then $Z_{p,N}^+(\varepsilon) = \tfrac{1}{2}(Z_p^\infty + R_{p,N}^\infty) < 0$
for all $p\leq\kappa$, i.e.\ $\mathrm{Re}(Z_p)<0$,
and consequently
\[
B_{\mathrm{int}}
= \sum_{p\leq\kappa}(\log p)^2 Z_{p,N}^+
 = \tfrac{1}{2}\sum_{p\leq\kappa}(\log p)^2
   \bigl(Z_p^\infty + R_{p,N}^\infty\bigr)
< 0.
\]
\end{proposition}

\begin{proof}
By~(A), $Z_p^\infty<0$. By~(B), $|R_{p,N}^\infty|<|Z_p^\infty|$.
Hence $Z_{p,N}^+ = \tfrac{1}{2}(Z_p^\infty + R_{p,N}^\infty) < 0$,
i.e.\ $\mathrm{Re}(Z_p)<0$.
Since $(\log p)^2>0$, the sum $B_{\mathrm{int}} < 0$.
In the infinite positive-ordinate limit, assumption~(A) gives directly
\[
B_{\mathrm{int}}^{\infty,+}
:=\tfrac{1}{2}\sum_{p\leq\kappa}(\log p)^2 Z_p^\infty(\varepsilon)<0,
\]
since $\kappa$ is finite and each weight $(\log p)^2$ is positive.
\end{proof}

\begin{remark}
The curvature quantity $B=\sum_{p\leq\kappa}(\log p)^2\mathrm{Re}(Z_p^{(2)})$
(Proposition~\ref{prop:decomp}) carries the additional $\gamma_k^2$~weights.
Since $\mathrm{Re}(Z_p)<0$ does not in general imply
$\mathrm{Re}(Z_p^{(2)})<0$, the Bias Conjecture alone does not give $B<0$.
The stronger statement $B=-19342.5<0$ is confirmed numerically at reference
parameters and constitutes independent evidence for the Curvature Bias.
\end{remark}

Assumption~(A) is the Bias Conjecture (Open Problem~\ref{op:bias}).
Assumption~(B) requires quantitative truncation-error control via
the bound of Definition~\ref{def:trunc}, which is not established here.
Neither condition is verified in this paper.

The logical chain from Conjecture~\ref{conj:bias} to the Curvature Bias
at $\sigma=\tfrac{1}{2}$ proceeds as follows (all conditional).
Under the Bias Conjecture and the truncation control of~(B),
Proposition~\ref{prop:sign} gives $B_{\mathrm{int}}<0$.
The full curvature quantity $B<0$ is confirmed numerically
($B=-19342.5$ at reference parameters).
Note that $B_{\mathrm{int}}<0$ does not imply $B<0$:
$B = \sum_p (\log p)^2\mathrm{Re}(Z_p^{(2)})$ carries additional
$\gamma_k^2$-weights absent from $B_{\mathrm{int}}$.
The missing link is Open Problem~\ref{op:wbridge} below.
By Reduction Lemma~\ref{lem:reduction}, $B<0$ establishes the
negative curvature bias.
Complementing this with analytic control of the stationarity condition
$|O'(\tfrac{1}{2})|\ll A(\varepsilon,N)\cdot\pi(\kappa)$ (Open Problem~\ref{op:stat})
would then complete the $\sigma=\tfrac{1}{2}$ Curvature Bias argument,
subject to all stated auxiliary conditions.

% -------------------------------------------------------
\section{Open Problems}
\label{sec:open}

We list five open problems in decreasing order of centrality
to the program of this series.

\begin{openproblem}[Bias Conjecture]\label{op:bias}
For each fixed prime $p$, prove that
$Z_p^\infty(\varepsilon)<0$
for all sufficiently small $\varepsilon>0$, where
$Z_p^\infty(\varepsilon)
 = 2\sum_{k=1}^\infty e^{-\varepsilon^2\gamma_k^2}\cos(\gamma_k\log p)$.
This is the central open problem of this paper.
Its proof via quantitative estimates on the Guinand--Weil explicit
formula --- controlling $\Gamma$-terms and correction terms relative
to the candidate main term $-(\log p)/(2\sqrt{\pi}\,\varepsilon\sqrt{p})$,
together with the GW bridge of Open Problem~\ref{op:gwbridge} if the
proof proceeds through the full Guinand--Weil object rather than
directly through the ordinate proxy $Z_p^\infty$,
--- would, in conjunction with the truncation control of
Proposition~\ref{prop:sign}(B), the Weighted Bias Bridge of
Open Problem~\ref{op:wbridge} (or an independent analytic proof
that $B<0$), and the Stationarity condition below,
complete the $\sigma=\tfrac{1}{2}$ Curvature Bias argument.
\end{openproblem}

\begin{openproblem}[Stationarity]\label{op:stat}
Fix $\varepsilon>0$ and $N\ge1$, and set
$A(\varepsilon,N)=\sum_{k=1}^N e^{-\varepsilon^2\gamma_k^2}$
as in Theorem~\ref{thm:trace}.
Establish whether there exists a constant $C_{\varepsilon,N}$,
independent of $\kappa$, such that
\[
  |O'(\tfrac12;\kappa,\varepsilon,N)|
  \le C_{\varepsilon,N}\,A(\varepsilon,N)\,\pi(\kappa)
\]
for all prime cutoffs $\kappa$.
At the reference parameters $(\kappa,\varepsilon,N)=(53,0.05,100)$,
$O'(\tfrac12)=+2.475$, approximately $11\%$ of $A(\varepsilon,N)\pi(\kappa)$
[Numerical].
An analytic proof would complement the Bias Conjecture.
The quantitative bounds of Carneiro, Chandee, and
Milinovich~\cite{CarneiroEtal13} on $S(t)$ provide a natural
analytic framework for future control of this bound.
\end{openproblem}

\begin{openproblem}[Weighted Bias Bridge]\label{op:wbridge}
Define $Z_{p,2}^\infty(\varepsilon)
:= 2\sum_{k=1}^\infty \gamma_k^2 e^{-\varepsilon^2\gamma_k^2}
   \cos(\gamma_k\log p)$.
Determine whether the sign implication
\[
  Z_p^\infty(\varepsilon)<0
  \quad\Longrightarrow\quad
  Z_{p,2}^\infty(\varepsilon)<0
\]
holds for sufficiently small $\varepsilon>0$ at fixed prime~$p$.
If the implication fails in general, identify additional hypotheses
under which the weighted sign transfer holds, or construct a
counterexample mechanism showing that the unweighted and
$\gamma_k^2$-weighted signs may decouple.
This ``weighted bias bridge'' is the missing link between the Bias
Conjecture (which controls $B_{\mathrm{int}}$) and the curvature
quantity $B$ (which carries $\gamma_k^2$-weights and enters
$O''(\tfrac{1}{2})=-2B$).
\end{openproblem}

\begin{openproblem}[Weighted spectral measure]\label{op:measure}
Assume $\|a_p\|>0$ for every active prime $p\leq\kappa$
(numerically satisfied throughout the tested range);
equivalently, restrict to the support
$S_\kappa^*:=\{p\leq\kappa:\|a_p\|>0\}$.
Let $K_\kappa:=D_a^{-1/2}G^{\mathrm{un}}D_a^{-1/2}$ and
$v_\kappa:=D_a^{1/2}c/\sqrt{\Estr(\kappa)}$.
Then $\|v_\kappa\|=1$ and
$\etaren(\kappa)=1-\langle v_\kappa,K_\kappa v_\kappa\rangle$ exactly.
Let $K_\kappa=\sum_j\lambda_j u_j u_j^T$; define
$\nu_\kappa^{(w)}:=\sum_j|\langle v_\kappa,u_j\rangle|^2\delta_{\lambda_j}$.
Then $\etaren(\kappa)=1-m_1(\nu_\kappa^{(w)})$ exactly.
The open problem is: establish tightness and convergence
of $\nu_\kappa^{(w)}$ as $\kappa\to\infty$ and identify
the limiting first moment.
\end{openproblem}

\begin{openproblem}[GW bridge]\label{op:gwbridge}
Establish whether the ordinate proxy
$Z_p^\infty(\varepsilon)=2\sum_{k=1}^\infty e^{-\varepsilon^2\gamma_k^2}\cos(\gamma_k\log p)$
agrees with the full Guinand--Weil object
$\widetilde Z_p^{\mathrm{GW}}(\varepsilon)$ to leading order
without assuming that all non-trivial zeros
lie on the critical line.
\end{openproblem}

% -------------------------------------------------------
\section*{Non-Circularity}
\label{sec:nocirc}

We document explicitly that the results of this paper do not rely
on the objects they are directed toward.

\textbf{No RH as input.}
At no point is the Riemann Hypothesis assumed.
The ordinates $\gamma_k$ are used as numerically specified inputs.
No property of their distribution requiring RH is assumed or used.

\textbf{No GUE, Montgomery, or Hilbert--P\'olya hypothesis.}
The operator $\Tt(\sigma)$ is constructed explicitly from the coupling
map $\Phi(\sigma)$ (Definition~\ref{def:coupling}).
No random matrix hypothesis, no GUE pair-correlation conjecture,
and no Hilbert--P\'olya postulate enters the construction or the proofs.

\textbf{$\Tt(\sigma)$ is not postulated.}
It is derived as $\Phi(\sigma)\circ\Phi(\sigma)^*$ from an explicit matrix.
It is not a hypothetical object.

\textbf{$\Tt$ is not a Hilbert--P\'olya operator.}
The numerical spectral diagnostic in \S\,\ref{ssec:spectral}
observes $\mu_j\approx\CT/\gamma_{k(j)}$ on the tested grid, hence
gives no support for a Hilbert--P\'olya interpretation; in any case
$\Tt(\sigma)$ is a finite-rank Gram-type operator built from external
ordinate inputs, not an operator postulated to have spectrum
$\{\gamma_k\}$.

\textbf{Theorem~\ref{thm:trace} is purely algebraic.}
The trace formula is proved by a two-line trigonometric manipulation.
No asymptotic analysis and no analytic number theory are used.

\textbf{The Bias Conjecture is open.}
Conjecture~\ref{conj:bias} is clearly marked as an open problem.
No theorem in this paper implies it.

\textbf{$\gamma_k$ as inputs, not conclusions.}
The ordinates appear as numerically specified inputs to the model.
The claim $\mathrm{Re}(\rho_k)=\tfrac{1}{2}$ is not used anywhere;
the reference point $\sigma_0=\tfrac{1}{2}$ is chosen as the object
of study, not assumed to be the unique zero location.

% -------------------------------------------------------
\section*{Acknowledgments}

The author thanks external reviewers and
independent verification checks.
Verification scripts are publicly available~\cite{Code}.

% -------------------------------------------------------
\section*{Call for Feedback}

This paper is part of an ongoing open research initiative toward the
Riemann Hypothesis.
All papers are freely available on Zenodo under CC~BY~4.0,
and the accompanying verification code is hosted on GitHub.

The author welcomes critical feedback, corrections, and collaboration.
All papers are available on Zenodo
(\url{https://zenodo.org/search?q=Ulrich+Tehrani}),
and the verification code on GitHub
(\url{https://github.com/utehrani/analysislab-nt}).

Topics of particular interest include: analytical approaches to the
Bias Conjecture via Guinand--Weil estimates; connections to existing
results on exponential sums over primes and zeros; and any errors or
gaps in the present arguments.

% -------------------------------------------------------
\newpage
\begin{thebibliography}{99}

\bibitem{Paper1}
U.~Tehrani,
``A Curvature Decomposition of the Explicit Formula
for the Riemann Zeta Function'',
Zenodo, \url{https://doi.org/10.5281/zenodo.19025598}, 2026.

\bibitem{Paper2}
U.~Tehrani,
``From Local Curvature to the Weil Functional:
An Explicit Construction'',
Zenodo, \url{https://doi.org/10.5281/zenodo.19106992}, 2026.

\bibitem{Paper3}
U.~Tehrani,
``A Finite-Cutoff Hilbert-Space Model
for Prime--Zero Energy Structure'',
Zenodo, \url{https://doi.org/10.5281/zenodo.19307989}, 2026.

\bibitem{Paper4}
U.~Tehrani,
``A Dual Operator for Prime--Zero Coupling
and a Conditional Proof of Energy Asymmetry'',
Zenodo, \url{https://doi.org/10.5281/zenodo.19364703}, 2026.

\bibitem{Code}
U.~Tehrani,
\emph{Verification scripts for the series},
GitHub, \url{https://github.com/utehrani/analysislab-nt}.

\bibitem{Guinand}
A.~P.~Guinand,
``Fourier reciprocities and the Riemann zeta-function'',
\emph{Proc.\ London Math.\ Soc.}\ (2) \textbf{51} (1950), 401--414.

\bibitem{Weil}
A.~Weil,
``Sur les `formules explicites' de la th\'eorie des nombres premiers'',
\emph{Comm.\ S\'em.\ Math.\ Univ.\ Lund, Suppl.}\ (1952), 252--265.

\bibitem{Mazhouda}
K.~Mazhouda,
``New asymptotic formulas for sums over zeros of functions
from the Selberg class'',
\emph{Ann.\ Sci.\ Math.\ Qu\'ebec} \textbf{34}(1) (2010), 85--108.
Corrected and revised version: arXiv:1505.01544v2 [math.NT], 2019.
Corrigendum: \emph{Ann.\ Math.\ Qu\'ebec}
\textbf{44} (2020), 197--200.
DOI: 10.1007/s40316-019-00124-3.

\bibitem{Hadamard}
J.~Hadamard,
``Sur la distribution des z\'eros de la fonction $\zeta(s)$
et ses cons\'equences arithm\'etiques'',
\emph{Bull.\ Soc.\ Math.\ France} \textbf{24} (1896), 199--220.

\bibitem{dlVP}
C.-J.~de la Vall\'ee Poussin,
``Recherches analytiques sur la th\'eorie des nombres premiers'',
\emph{Ann.\ Soc.\ Sci.\ Bruxelles} \textbf{20} (1896), 183--256.

\bibitem{Connes}
A.~Connes,
``Trace formula in noncommutative geometry and the zeros
of the Riemann zeta function'',
\emph{Selecta Math.\ (N.S.)} \textbf{5}(1) (1999), 29--106.

\bibitem{Burnol}
J.-F.~Burnol,
``On Fourier and Zeta(s)'',
\emph{Forum Math.} \textbf{16}(6) (2004), 789--840.

\bibitem{RubinsteinSarnak}
M.~Rubinstein and P.~Sarnak,
``Chebyshev's bias'',
\emph{Experiment.\ Math.} \textbf{3}(3) (1994), 173--197.

\bibitem{ConnesConsani21}
A.~Connes and C.~Consani,
``Weil positivity and trace formula, the archimedean place,''
\textit{Selecta Math.\ (N.S.)} \textbf{27} (2021), no.~77.
arXiv:2006.13771.
DOI: 10.1007/s00029-021-00689-4.

\bibitem{ConnesMoscovici22}
A.~Connes and H.~Moscovici,
``The UV prolate spectrum matches the zeros of zeta,''
\textit{Proc.\ Natl.\ Acad.\ Sci.\ USA}
\textbf{119}, no.~22 (2022).
DOI: 10.1073/pnas.2123174119.

\bibitem{CarneiroEtal13}
E.~Carneiro, V.~Chandee, and M.~B.~Milinovich,
``Bounding $S(t)$ and $S_1(t)$ on the Riemann hypothesis,''
\textit{Math.\ Ann.} \textbf{356} (2013), 939--968.
arXiv:1309.1526.
DOI: 10.1007/s00208-012-0876-z.

\bibitem{CCM25}
A.~Connes, C.~Consani, and H.~Moscovici,
``Zeta spectral triples,''
arXiv:2511.22755 (2025).

\bibitem{Yakaboylu26}
E.~Yakaboylu,
``Nontrivial Riemann Zeros as Spectrum,''
arXiv:2408.15135, v15 (March 2026, preprint).

\end{thebibliography}

% -------------------------------------------------------
\enlargethispage{3\baselineskip}
\vspace*{\fill}
\begin{minipage}{\linewidth}
\rule{\linewidth}{0.4pt}\smallskip\\
\small\emph{Paper~5 v1.19 \textperiodcentered\ May~2026
\textperiodcentered\ Pauca sed matura}
\end{minipage}

\end{document}
